% !TeX spellcheck = en_GB
\chapter{Environments}
\label{chap:environment}

\section{Reinforcement learning}
Reinforcement learning might be best introduced by description of the problem it aims to solve.
\begin{defn}[Markov’s decision problem]
We define Markov’s decision problem (MDP) by a four-tuple (S, A, P, R) where:
\begin{itemize}
	\item S is a set of all possible states of the system
	\item A is a set of all possible actions available in the system
	\item P is a transition function,$ P(s, a,  s')$ being probability of transition from s to s' give action a is executed.
	\item R is reward function, $ R(s, a, s') $ being the reward received when the execution of action a in state s leads to state s'.
\end{itemize}
\end{defn}
To tie this into the information from chapter \ref{ch:agents} $ A=A(E)$, $ R $ is performance measure of an agent solving the MDP and the percepts correspond to real states of the environment. This would not be the case for partially observable MDP.
\begin{defn}
POMDP
\end{defn}
\begin{defn}
Best policy
\end{defn}
The goal of reinforcement learning is then to proximate best policy by trial and error, based on the received reward, as accurately as possible. 
This can be done utilising various methods such as Q-learning, but importantly for this thesis it could be  



\section{Gymnasium}
In this thesis we will set our agents into environments provided by Gymnasium project which emerged from it's predecessor OpenAI Gym project. Both of them Python libraries, trying to provide researches with benchmark reinforcement learning problems in form of simple games, sometimes even recreating old console games such as Air Raid or Backgammon.

While Gymnasium provides wide range of environments, quite a number of them if not most do not provide their outputs in structured form. Instead the resulting screen has to be processed by user defined visual processor to extract information crucial for completing the task. 

This introduces a layer of complexity that is outside of the scope of this thesis. Therefore we have decided to narrow our selection only to environments which do provide structured output. This is the two environments we have eventually selected for our study:

\subsection{Cartpole}
\label{env:cartpole}

Cartpole is one of the most used benchmark environments in the Gymnasium library. It's a simple environment where the agent controls a cart rolling on flat surface with a pole on top.
The goal is to keep the pole balanced on top of the cart. The reward is calculated every step depending on the angle of the stick. This means we have bounded 
\begin{defn}
	An agent that achieves score 500 in Cartpole is called solution or solving individual.
\end{defn}
\subsubsection{Observation space}
\subsubsection{Action space}
\subsection{Lunar Lander}
The other environment we have selected is Lunar Lander where the agent controls descend of a landing shuttle on rugged planetary (or the Moon's) surface with the help of strangely placed rockets. The points are subtracted for staying in the air or for getting the shuttle into a spin. Bonus points awarded for landing in the center of the  The shuttle has infinite fuel so the game ends only when the shuttle lands safely or crushes into the ground. Luckily we can at least limit the number of steps for which the environment is allowed to go on. For a neural network training this however poses a specific kind of problem. Firstly very narrow behavioral window that can achieve positive fitness. But mainly as opposed to the Cartpole environment (\ref{env:cartpole}), now the unsuccesful runs may take significantly longer time than the successful ones. This means experiments are particularly costly for slowly converging algorithms. The upper limit of fitness is 200, the lower would possibly be $-\inf$, depending on the lenght limitation imposed.
We therefore differntiate between two stages of obtaining successful solution.
\begin{defn}
	An agent that achieves score 200 in Lunar lander is called solution or solving individual.
\end{defn}
\begin{defn}
	\label{def:wellbehaved}
	An agent that achieves positive score  in Lunar lander is called well-behaved.
\end{defn}
 Definition \ref{def:wellbehaved} stems from the fact that it can only be achieved if the agent does not crash the shuttle as doing so yields a penalisation of -100 points while not obtaining a boost of +100 points for safe landing, driving even a possible solution with score 200 to 0.
\subsubsection{Observation space}
\subsubsection{Action space}




 
 
